% Enterprise RAG Assistant SRS
% Compile main.tex with pdfLaTeX twice, or use latexmk -pdf main.tex.
% Content edition 1.1 dated 28 October 2026.
\documentclass[11pt,letterpaper]{article}
\usepackage[T1]{fontenc}
\usepackage[utf8]{inputenc}
\usepackage{mathptmx}
\usepackage[margin=0.8in,headheight=14pt]{geometry}
\usepackage{array,longtable,booktabs}
\usepackage[table]{xcolor}
\usepackage{enumitem}
\usepackage{titlesec}
\usepackage{fancyhdr}
\usepackage{needspace}
\usepackage{microtype}
\usepackage{xurl}
\usepackage[hidelinks,bookmarksnumbered,pdfauthor={Abdul Samad},pdftitle={Enterprise RAG Assistant Software Requirements Specification}]{hyperref}
\usepackage{lastpage}
\setlength{\parindent}{0pt}
\setlength{\parskip}{5pt}
\linespread{1.05}
\setlength{\emergencystretch}{2em}
\setlength{\tabcolsep}{5pt}
\renewcommand{\arraystretch}{1.13}
\setlength{\LTpre}{6pt}
\setlength{\LTpost}{7pt}
\arrayrulecolor{black!20}
\setlist[itemize]{leftmargin=1.5em,itemsep=2pt,topsep=3pt}
\setcounter{secnumdepth}{3}
\setcounter{tocdepth}{2}
\titleformat{\section}{\large\bfseries}{\thesection}{0.6em}{}
\titleformat{\subsection}{\normalsize\bfseries}{\thesubsection}{0.6em}{}
\titleformat{\subsubsection}{\normalsize\bfseries}{\thesubsubsection}{0.6em}{}
\titlespacing*{\section}{0pt}{12pt}{6pt}
\titlespacing*{\subsection}{0pt}{10pt}{5pt}
\titlespacing*{\subsubsection}{0pt}{9pt}{4pt}
\pagestyle{fancy}
\fancyhf{}
\fancyhead[L]{\footnotesize Enterprise RAG Assistant}
\fancyhead[R]{\footnotesize Software Requirements Specification}
\fancyfoot[R]{\footnotesize Version 1.1\quad Page \thepage\ of \pageref{LastPage}}
\renewcommand{\headrulewidth}{0pt}
\renewcommand{\footrulewidth}{0pt}
\newcolumntype{P}[1]{>{\raggedright\arraybackslash}p{#1}}
\begin{document}
\begin{titlepage}
\centering
{\bfseries Institute of Management Sciences, Peshawar\par}
\vspace{6pt}
BS Software Engineering\par
Software Requirements Engineering Lab\par
\vspace{1.0cm}
{\LARGE\bfseries Enterprise RAG Assistant\par}
\vspace{9pt}
{\Large Software Requirements Specification\par}
\vspace{9pt}
Version 1.1\par
Complete planned version 1\par
\vspace{1.0cm}
Prepared by\par
{\bfseries Abdul Samad\par}
Student ID 243203624\par
\vspace{9pt}
Submitted to\par
{\bfseries Sir Fahim Muhammad Khan\par}
28 October 2026\par
\vspace{1.0cm}
Status: Prepared for review and approval\par
\end{titlepage}
\setcounter{page}{2}
\section*{Document Control}

\addcontentsline{toc}{section}{Document Control}

This SRS defines the planned first release of Enterprise RAG Assistant and its acceptance criteria. Document approval establishes the requirements baseline; product acceptance follows implementation and verification.\par

\begingroup\small

\begin{longtable}{|P{\dimexpr0.2279\linewidth-0.9118\tabcolsep-0.6838\arrayrulewidth\relax}|P{\dimexpr0.7721\linewidth-3.0882\tabcolsep-2.3162\arrayrulewidth\relax}|}

\hline\rowcolor{black!8}

\textbf{Item} & \textbf{Value} \\ \hline

\endfirsthead

\hline\rowcolor{black!8}

\textbf{Item} & \textbf{Value} \\ \hline

\endhead

\hline\endfoot

\hline\endlastfoot

Document identifier & ERA-SRS-001 \\ \hline

Product & Enterprise RAG Assistant \\ \hline

Version and date & 1.1 | 28 October 2026 \\ \hline

Author and owner & Abdul Samad | Student ID 243203624 \\ \hline

Institution & Institute of Management Sciences, Peshawar \\ \hline

Programme and course & BS Software Engineering | Software Requirements Engineering Lab \\ \hline

Reviewer & Sir Fahim Muhammad Khan \\ \hline

Approval status & Pending the records in Section 5 \\ \hline

Scope baseline & Complete planned version 1; one organisation and four simulated employee roles \\ \hline

\end{longtable}

\endgroup

\subsection*{Revision History}

\begingroup\small

\begin{longtable}{|P{\dimexpr0.0956\linewidth-0.7647\tabcolsep-0.4779\arrayrulewidth\relax}|P{\dimexpr0.1544\linewidth-1.2353\tabcolsep-0.7721\arrayrulewidth\relax}|P{\dimexpr0.5588\linewidth-4.4706\tabcolsep-2.7941\arrayrulewidth\relax}|P{\dimexpr0.1912\linewidth-1.5294\tabcolsep-0.9559\arrayrulewidth\relax}|}

\hline\rowcolor{black!8}

\textbf{Version} & \textbf{Date} & \textbf{Description} & \textbf{Author} \\ \hline

\endfirsthead

\hline\rowcolor{black!8}

\textbf{Version} & \textbf{Date} & \textbf{Description} & \textbf{Author} \\ \hline

\endhead

\hline\endfoot

\hline\endlastfoot

1.0 & 27 October 2026 & Initial formal SRS assembled from the project requirements v0.2, design v0.3, data-source records, and the supplied course guidance. & Abdul Samad \\ \hline

1.1 & 28 October 2026 & Shortened repeated context and verification material; retained baseline requirements, acceptance criteria and open decisions. Added an equivalent LaTeX edition. & Abdul Samad \\ \hline

\end{longtable}

\endgroup

\textbf{Interpretation.} Shall denotes a requirement, should a desired behaviour, and may an option. Provisional targets require the measurement conditions in Section 6.3 before acceptance.\par

\textbf{Requirement control.} FR-01 to FR-13 and NFR-01 to NFR-07 retain their existing identifiers. BR, IR, DR, and DC identify supporting business, interface, data, and operating rules.\par

\textbf{Priority.} High marks a version 1 obligation. Medium marks a provisional quality target. These priorities are subject to review.\par

\clearpage

\tableofcontents

\clearpage

\clearpage

\section{Introduction}

\subsection{Purpose}

Enterprise RAG Assistant answers employee questions from approved company documents. This SRS defines its required behaviour, permissions, interfaces, data rules, and quality targets. It gives the owner, developer, tester, and course reviewer a common basis for implementation and acceptance.\par

The owner uses the SRS to control scope. Developers use Section 3 to implement the system; testers and reviewers use Section 6 to check coverage and record results.\par

\subsection{Scope}

The assistant accepts questions in English and returns answers supported by documents the employee may access. A substantive answer includes a source citation. If the available evidence is insufficient, the system returns a generic not-found result.\par

Version 1 covers offline Markdown and text-based PDF ingestion, department tagging, signed-token login, restricted question answering, a REST API, browser chat, query metrics, evaluation, and a deployed demo. Acme Corp is a fictional company represented by selected public GitLab Handbook pages.\par

Voice, multilingual support, native mobile apps, model training, document editing through chat, live company-system synchronisation, SSO, and multi-tenant hosting are excluded. OCR for scanned PDFs is a possible later extension.\par

\subsection{Definitions Acronyms and Abbreviations}

\begingroup\small

\begin{longtable}{|P{\dimexpr0.1985\linewidth-0.7941\tabcolsep-0.5956\arrayrulewidth\relax}|P{\dimexpr0.8015\linewidth-3.2059\tabcolsep-2.4044\arrayrulewidth\relax}|}

\hline\rowcolor{black!8}

\textbf{Term} & \textbf{Meaning in this specification} \\ \hline

\endfirsthead

\hline\rowcolor{black!8}

\textbf{Term} & \textbf{Meaning in this specification} \\ \hline

\endhead

\hline\endfoot

\hline\endlastfoot

SRS & Software Requirements Specification; the controlled statement of required product behaviour. \\ \hline

RAG & Retrieval-augmented generation; answering a question using document passages retrieved for that question. \\ \hline

LLM & Large language model used to produce an answer from the supplied evidence. \\ \hline

RBAC & Role-based access control; permission determined from an authenticated user role. \\ \hline

Document and chunk & A source file and a smaller passage derived from that file. A chunk retains its document identity and department. \\ \hline

Citation & A reference to an authorised source, including its document title and section or PDF page. \\ \hline

Not found & The same generic outcome for insufficient evidence, restricted material, or non-existent material; it does not reveal whether a restricted document exists. \\ \hline

Faithfulness & The degree to which an answer is supported by its supplied source passages, as scored by the agreed evaluation method. \\ \hline

API and bearer token & A software interface and the signed access token included with authenticated requests. \\ \hline

Checksum & A SHA-256 fingerprint used to identify a source document revision. \\ \hline

BM25 and reranker & Keyword-based retrieval scoring and a later step that reorders candidate passages for relevance. \\ \hline

RRF & Reciprocal rank fusion; a method of merging ranked retrieval results. \\ \hline

FR and NFR & Functional requirement and non-functional requirement. \\ \hline

BR IR DR DC & Business rule, interface rule, data rule, and design or operating constraint. \\ \hline

TBD & To be decided and recorded before the affected acceptance criterion is approved. \\ \hline

\end{longtable}

\endgroup

\subsection{References}

References R1 and R2 supply the course structure; R3 to R6 establish the project baseline. R1 discusses ISO/IEC/IEEE 29148, and R2 identifies its outline as adapted from IEEE Std 830-1993. This SRS follows the supplied guidance without claiming certified conformity.\par

\begingroup\small

\begin{longtable}{|P{\dimexpr0.0662\linewidth-0.3971\tabcolsep-0.2647\arrayrulewidth\relax}|P{\dimexpr0.6471\linewidth-3.8824\tabcolsep-2.5882\arrayrulewidth\relax}|P{\dimexpr0.2868\linewidth-1.7206\tabcolsep-1.1471\arrayrulewidth\relax}|}

\hline\rowcolor{black!8}

\textbf{Ref} & \textbf{Document and source} & \textbf{Version or date} \\ \hline

\endfirsthead

\hline\rowcolor{black!8}

\textbf{Ref} & \textbf{Document and source} & \textbf{Version or date} \\ \hline

\endhead

\hline\endfoot

\hline\endlastfoot

R1 & Software Requirement Specification (SRS); supplied course handout, SRS.pdf. & Four pages; date not stated. \\ \hline

R2 & Software Requirements Specification (SRS) Template; CS3911, supplied PDF. & 30 pages; header 08/15/22. \\ \hline

R3 & Requirements Document: Enterprise RAG Assistant; Abdul Samad; docs/01-requirements.md. & Version 0.2. \\ \hline

R4 & Design Document: Enterprise RAG Assistant; Abdul Samad; docs/02-design.md. & Version 0.3. \\ \hline

R5 & Data Sources and manifest; docs/data-sources.md and docs/data-manifest.csv. & Dataset downloaded 11 October 2026. \\ \hline

R6 & Project README, AGENTS.md, application source, and tests. & Reviewed 27 October 2026. \\ \hline

R7 & GitLab Handbook; GitLab; source revision and locations in Section 6.4. & Pinned commit in Section 6.4. \\ \hline

R8 & Handbook attribution record in the project README and R5. & CC BY-SA 4.0; URL in Section 6.4. \\ \hline

\end{longtable}

\endgroup

Project repository: \url{https://github.com/abdulsamad-codes/enterprise-rag-assistant}. References R3 to R6 are in the project repository; the supplied PDFs accompany this document.\par

\subsection{Overview}

Section 2 sets the product context and release boundary. Section 3 specifies behaviour and quality. Sections 4 and 5 govern changes and approval. Section 6 holds verification, traceability, provenance, and unresolved decisions.\par

\clearpage

\section{The Overall Description}

\subsection{Product Perspective}

The assistant serves one organisation and supplements its document collection. It does not edit policies or perform the payroll, leave, purchasing, or other transactions described in them. Employees follow the cited policy and its existing approval process.\par

The system includes ingestion, authentication, access filtering, retrieval, generation, citation checking, the API, chat, logging, and evaluation. It reads approved files offline. At query time, the external model service receives the question and authorised evidence; the assistant does not browse the live handbook.\par

\begingroup\small

\begin{longtable}{|P{\dimexpr0.1985\linewidth-1.1912\tabcolsep-0.7941\arrayrulewidth\relax}|P{\dimexpr0.3897\linewidth-2.3382\tabcolsep-1.5588\arrayrulewidth\relax}|P{\dimexpr0.4118\linewidth-2.4706\tabcolsep-1.6471\arrayrulewidth\relax}|}

\hline\rowcolor{black!8}

\textbf{External actor or resource} & \textbf{Input to the assistant} & \textbf{Output from the assistant} \\ \hline

\endfirsthead

\hline\rowcolor{black!8}

\textbf{External actor or resource} & \textbf{Input to the assistant} & \textbf{Output from the assistant} \\ \hline

\endhead

\hline\endfoot

\hline\endlastfoot

Employee & Credentials; English question; signed token on later requests. & Login result; answer and citations; generic not-found result. \\ \hline

Project operator & Approved documents; department mappings; deployment settings; evaluation set. & Ingestion result; health status; query metrics; evaluation report. \\ \hline

Language-model service & Generated answer and candidate citation identifiers returned to the assistant. & Question, generation instructions, and only authorised evidence sent to the service. \\ \hline

Approved document collection & Markdown files and text-based PDF files. & No source-document edits are produced by the assistant. \\ \hline

\end{longtable}

\endgroup

The table describes the external context. Internal retrieval and storage components are specified in R4.\par

\subsubsection{System Interfaces}

The employee interface uses the API for login and questions. Ingestion reads local files; evaluation exercises the system under recorded roles. Version 1 has no payroll, HR-system, SSO-provider, or live document-management integration.\par

\subsubsection{User Interfaces}

Employees use login and chat screens showing questions, answers, citations, and request outcomes. They cannot assign themselves another role or add files to the approved collection. Operators use command-line ingestion and evaluation tools.\par

\subsubsection{Hardware Interfaces}

No specialised hardware interfaces are required. The application uses ordinary storage, networking, and a browser; it does not control external devices.\par

\subsubsection{Software Interfaces}

The backend exchanges requests with a hosted model service. R4 selects Groq behind a replaceable client interface. Storage engines, retrieval libraries, and the UI framework are internal design choices. Provider versions and limits are deployment settings subject to OD-03 and OD-04.\par

\subsubsection{Communications Interfaces}

The deployed client and API communicate over HTTPS; authenticated requests carry a bearer token. The backend needs outbound access to the model provider. No custom device or messaging protocol is used.\par

\subsubsection{Memory Constraints}

RAM, storage, throughput, and concurrent-user limits are not yet baselined. They will be measured against the selected models, document collection, and host under OD-01.\par

\subsubsection{Operations}

The operator prepares documents offline while serving is stopped. After ingestion, the service restarts and checks consistency between the chunk collection and retrieval indexes. Employees then log in and ask questions. Evaluation is a separate operator-run activity.\par

Retention, backup, and recovery arrangements remain open under OD-02. Continuous service during document updates is outside the first-release commitment.\par

\subsubsection{Site Adaptation Requirements}

The operator configures source paths, department mapping, accounts, provider credentials, storage, retrieval settings, and content logging. Private company deployment also requires an agreed access map and approval of the provider data flow. No physical site changes are required.\par

\subsection{Product Functions}

\begin{itemize}

\item Prepare an approved document collection and preserve its department and source metadata.

\item Authenticate employees and determine their permissions from a backend-verified token.

\item Find relevant passages within the departments permitted for the employee.

\item Produce an evidence-based answer with validated citations, or return a generic not-found outcome.

\item Expose the behaviour through an API and chat interface.

\item Record operational metrics and compare answer quality across retrieval configurations.

\end{itemize}

\subsection{User Characteristics}

Employees need English reading skills and ordinary browser experience, not knowledge of retrieval or language models. The roles below determine permissions, not technical skill. The operator must be able to configure the deployment and run its command-line tools.\par

\begingroup\small

\begin{longtable}{|P{\dimexpr0.1176\linewidth-0.9412\tabcolsep-0.5882\arrayrulewidth\relax}|P{\dimexpr0.1471\linewidth-1.1765\tabcolsep-0.7353\arrayrulewidth\relax}|P{\dimexpr0.3824\linewidth-3.0588\tabcolsep-1.9118\arrayrulewidth\relax}|P{\dimexpr0.3529\linewidth-2.8235\tabcolsep-1.7647\arrayrulewidth\relax}|}

\hline\rowcolor{black!8}

\textbf{Role} & \textbf{Token value} & \textbf{Permitted departments} & \textbf{Restricted departments} \\ \hline

\endfirsthead

\hline\rowcolor{black!8}

\textbf{Role} & \textbf{Token value} & \textbf{Permitted departments} & \textbf{Restricted departments} \\ \hline

\endhead

\hline\endfoot

\hline\endlastfoot

New Hire & new\_hire & onboarding, general & people, finance, engineering \\ \hline

HR & hr & people, onboarding, general & finance, engineering \\ \hline

Finance & finance & finance, onboarding, general & people, engineering \\ \hline

Engineer & engineer & engineering, onboarding, general & people, finance \\ \hline

\end{longtable}

\endgroup

Operator is a maintenance responsibility, not a fifth employee role. The handbook is public; its access restrictions are simulated for this project.\par

\subsection{Constraints}

\begin{itemize}

\item The first release serves one organisation with four employee roles and department-level permissions.

\item Permission checks occur in the backend before evidence reaches answer generation. Client controls and prompts cannot grant access.

\item The serving process uses one API worker. Document ingestion is performed offline.

\item The demo collection is pinned to a recorded handbook commit. Dataset provenance and attribution are retained.

\item External service availability and account limits affect generation and evaluation.

\item Secrets, logs, local datasets, virtual environments, and generated databases are excluded from the public source repository.

\end{itemize}

\subsection{Assumptions and Dependencies}

The baseline comes from the owner's goals and project specification; no enterprise stakeholder interview has been completed. The operator is responsible for source accuracy, approval, and classification.\par

Employees need a browser and network access; generation depends on the model provider. Roles are trusted only after backend verification. A revoked account may retain access until its token expires; immediate revocation is deferred.\par

Quality targets remain provisional until the environment and scoring protocol are recorded. Changes to these assumptions follow Section 4.\par

\subsection{Apportioning of Requirements}

\begingroup\small

\begin{longtable}{|P{\dimexpr0.2059\linewidth-0.8235\tabcolsep-0.6176\arrayrulewidth\relax}|P{\dimexpr0.7941\linewidth-3.1765\tabcolsep-2.3824\arrayrulewidth\relax}|}

\hline\rowcolor{black!8}

\textbf{Release} & \textbf{Included capabilities} \\ \hline

\endfirsthead

\hline\rowcolor{black!8}

\textbf{Release} & \textbf{Included capabilities} \\ \hline

\endhead

\hline\endfoot

\hline\endlastfoot

Version 1 baseline & Approved local Markdown and text-based PDF ingestion; four simulated roles; signed-token login; restricted retrieval; grounded answers and citations; API and chat; metrics; evaluation; deployed demo. \\ \hline

Later versions & Enterprise SSO and immediate entitlement checks; finer access groups; background document synchronisation; multiple API workers; optional grounding verification at request time; automatic personal-data redaction; per-client configuration. \\ \hline

Outside the present commitment & Voice, multilingual support, native mobile applications, model fine-tuning, document editing through chat, multi-tenant hosting, and OCR for scanned PDFs. \\ \hline

\end{longtable}

\endgroup

Later-release capabilities require an approved scope change before becoming version 1 obligations.\par

\clearpage

\section{Specific Requirements}

This section specifies the complete planned version 1. All functional and supporting rules are High priority. NFR priorities are stated separately. Source references identify their basis; verification IDs name planned procedures, not completed tests.\par

\subsection{External Interface Requirements}

\subsubsection{User Interfaces}

\begingroup\small

\begin{longtable}{|P{\dimexpr0.0956\linewidth-0.5735\tabcolsep-0.3824\arrayrulewidth\relax}|P{\dimexpr0.7721\linewidth-4.6324\tabcolsep-3.0882\arrayrulewidth\relax}|P{\dimexpr0.1324\linewidth-0.7941\tabcolsep-0.5294\arrayrulewidth\relax}|}

\hline\rowcolor{black!8}

\textbf{ID} & \textbf{Required interface behaviour} & \textbf{Verification} \\ \hline

\endfirsthead

\hline\rowcolor{black!8}

\textbf{ID} & \textbf{Required interface behaviour} & \textbf{Verification} \\ \hline

\endhead

\hline\endfoot

\hline\endlastfoot

IR-01 & The employee interface shall provide credential entry, question entry, answer display, citation display, and visible feedback when authentication or a request fails. [R3 FR-07, FR-11; R4 Section 2] & AT-07, AT-11 \\ \hline

IR-02 & The employee interface shall send the backend-issued token with question requests and shall not use a selectable role value as authority for access. [R4 Section 4.2] & AT-05, AT-11 \\ \hline

IR-03 & The interface shall present the same not-found outcome for restricted and missing evidence. [R3 FR-12] & AT-12 \\ \hline

\end{longtable}

\endgroup

Browser support, screen assumptions, and accessibility scope remain OD-07. The required login, question, answer, and citation interactions apply regardless of the final layout.\par

\subsubsection{Hardware Interfaces}

No direct hardware interfaces apply. Host capacity is governed by OD-01.\par

\subsubsection{Software Interfaces}

\begingroup\small

\begin{longtable}{|P{\dimexpr0.1765\linewidth-1.0588\tabcolsep-0.7059\arrayrulewidth\relax}|P{\dimexpr0.3897\linewidth-2.3382\tabcolsep-1.5588\arrayrulewidth\relax}|P{\dimexpr0.4338\linewidth-2.6029\tabcolsep-1.7353\arrayrulewidth\relax}|}

\hline\rowcolor{black!8}

\textbf{Operation} & \textbf{Input and authority} & \textbf{Output and conditions} \\ \hline

\endfirsthead

\hline\rowcolor{black!8}

\textbf{Operation} & \textbf{Input and authority} & \textbf{Output and conditions} \\ \hline

\endhead

\hline\endfoot

\hline\endlastfoot

POST /auth/login & User credentials; no prior bearer token required. & Authentication result and a signed token on success. Credential field names and error schema remain TBD. \\ \hline

POST /ask & JSON body containing question; a verified bearer token supplies the role. & Answer, citations, and a found flag. The response distinguishes a request failure from an ordinary not-found outcome. \\ \hline

GET /health & No authentication required in the current design. & Service status and consistency status of the retrieval data; public response detail remains TBD. \\ \hline

Ingestion CLI & Configured data directory; department and provenance information. & Prepared document data and retrieval indexes, or an ingestion error. No remote ingestion endpoint is provided. \\ \hline

Evaluation CLI & Question records with id, role, question, expected\_answer, source\_doc, should\_refuse. & Per-question outcomes and aggregate evaluation results. \\ \hline

Model provider & Question, generation instructions, and authorised source passages with chunk identifiers. & Generated answer and proposed citation identifiers, or a provider error. \\ \hline

\end{longtable}

\endgroup

\textbf{IR-04.} The question API shall derive the role from the verified token rather than from the question body. [R3 FR-11; R4 Section 6] Verification: AT-11.\par

\textbf{IR-05.} The model request shall contain no source passage outside the authenticated role's permitted departments. [R3 FR-05; R4 Section 4.1] Verification: AT-05 with captured provider input.\par

The question field is text. Final JSON fields, credential format, input limits, error codes, and invalid-input behaviour remain OD-04. These must be fixed before API-contract approval.\par

\subsubsection{Communications Interfaces}

\textbf{IR-06.} The deployed user-facing service shall use HTTPS and carry its signed access token as a bearer token on authenticated requests. [R4 architecture and Section 4.2] Verification: AT-11 and deployment inspection.\par

Provider secrets belong in deployment configuration. Timeouts, retries, and service-error messages remain OD-04.\par

\subsection{System Features and Functional Requirements}

Each feature states its purpose and operating context, followed by requirements and acceptance criteria. Original IDs are retained for traceability to R3.\par

\subsubsection{Document Ingestion and Classification}

\textbf{Purpose.} Prepare approved documents while retaining source identity and department classification.\par

\textbf{Operating context.} The operator runs ingestion while serving is stopped. The folder must be readable and source paths mapped. Ingestion extracts text, prepares labelled chunks, updates changed documents, and removes deleted sources.\par

\Needspace{7\baselineskip}

\textbf{FR-01 Supported document ingestion}\par

The system shall ingest Markdown documents and text-based PDF documents from a configured local data folder.\par

\textbf{Acceptance.} AT-01: ingest one readable Markdown fixture and one text-based PDF fixture; confirm that their body text is available in the prepared collection and remains associated with the source file. Scanned PDFs are outside the required set.\par

\Needspace{7\baselineskip}

\textbf{FR-02 Department and source metadata}\par

The system shall associate each ingested document and its derived chunks with a department and source metadata that supports access enforcement and source citation.\par

\textbf{Acceptance.} AT-02: verify department, document identity, source path, checksum, title, and section or page metadata for prepared fixtures; reject a source that cannot be assigned an approved department.\par

\textbf{Exception handling.} Malformed, excluded, or unmapped paths receive no default department. Ingestion errors must not expose unclassified content. Corrupt-file handling and partial-run reporting remain OD-04.\par

\subsubsection{Authentication and Role Enforcement}

\textbf{Purpose.} Authenticate employees and enforce their permissions before evidence reaches generation.\par

\textbf{Operating context.} A valid demo account has an assigned role. Login issues a signed token; the backend verifies it on each question and selects permitted departments. Unknown roles receive no fallback access. The design sets a 30-minute expiry.\par

\Needspace{7\baselineskip}

\textbf{FR-11 Signed-token authentication}\par

The system shall authenticate an employee and determine the employee's role from a signed token verified by the backend.\par

\textbf{Acceptance.} AT-11: valid credentials produce a usable signed token; altered, invalid, and expired tokens cannot authorise a question. A role supplied in the question body cannot change the role determined by the token.\par

\Needspace{7\baselineskip}

\textbf{FR-05 Retrieval within permitted departments}\par

The system shall supply retrieval evidence only from documents within the authenticated employee's permitted departments.\par

\textbf{Acceptance.} AT-05: for each role, inspect retrieved chunks and the complete model request; every passage must belong to an allowed department. A search dominated by disallowed documents must still consider permitted candidates before selecting the final results.\par

\Needspace{7\baselineskip}

\textbf{FR-12 Indistinguishable restricted and missing outcomes}\par

The system shall return the same generic not-found reply when requested evidence is restricted and when it does not exist in the approved collection.\par

\textbf{Acceptance.} AT-12: ask equivalent questions about a known restricted source and a non-existent source as the same role; compare the public found flag, answer message, and citations. The not-found result must not disclose the restricted title or its existence.\par

\subsubsection{Question Answering and Evidence Sufficiency}

\textbf{Purpose.} Answer from permitted evidence and refuse unsupported conclusions.\par

\textbf{Operating context.} An authenticated employee submits a question to a ready service. Relevant permitted passages are assessed for sufficiency. Insufficient evidence produces not found; supported generation proceeds to citation checking.\par

\Needspace{7\baselineskip}

\textbf{FR-03 Answers grounded in retrieved documents}\par

The system shall answer a natural-language question using the retrieved document content as its source of factual evidence.\par

\textbf{Acceptance.} AT-03: submit factual and multi-section questions with expected answers derived from approved sources; assess answer correctness and source support using the evaluation method defined for NFR-02.\par

\Needspace{7\baselineskip}

\textbf{FR-06 Refusal when evidence is insufficient}\par

The system shall return a not-found outcome when the permitted retrieved evidence does not support an answer.\par

\textbf{Acceptance.} AT-06: below-threshold evidence stops generation; an insufficient-context result from generation also becomes not found. Questions outside the collection must not produce a factual answer based solely on model knowledge.\par

\Needspace{7\baselineskip}

\textbf{FR-10 Resistance to instruction overrides}\par

The system shall preserve its access rules and evidence-based answering behaviour when a question or retrieved document contains instructions that attempt to override them.\par

\textbf{Acceptance.} AT-10: test instructions embedded in questions and document passages that request another role, restricted sources, unsupported answers, or invented citations. Inspect both the public result and the model input; no restricted passage may be supplied.\par

\subsubsection{Citation Validation and Answer Delivery}

\textbf{Purpose.} Trace each substantive answer to evidence actually supplied for generation.\par

\textbf{Operating context.} The generator returns an answer and citation IDs. The system retains only IDs from the supplied chunks. If none remain, it suppresses the answer and returns not found without another model call to repair citations.\par

\Needspace{7\baselineskip}

\textbf{FR-04 Readable source citations}\par

The system shall include at least one citation identifying the document title and section or PDF page in every substantive answer returned to an employee.\par

\textbf{Acceptance.} AT-04: each successful found=true response includes a citation with a document title and section or page that resolves to supplied evidence. A not-found response is not a substantive answer and requires no invented citation.\par

\Needspace{7\baselineskip}

\textbf{FR-13 Citation identifier validation}\par

The system shall validate every proposed citation against the chunks supplied to the model, discard invalid identifiers, and return not found when no valid citation remains.\par

\textbf{Acceptance.} AT-13: all-valid identifiers are retained; mixed valid and invalid identifiers retain only valid entries; all-invalid or empty identifiers produce not found. Confirm that citation repair does not trigger another model call.\par

\Needspace{7\baselineskip}

\textbf{FR-07 API and chat access}\par

The system shall make question answering available through a REST API and a browser-based chat interface.\par

\textbf{Acceptance.} AT-07: authenticate and ask the same supported question through each interface; confirm equivalent access restrictions, answer outcome, and citation requirements. Confirm that ingestion and evaluation remain operator tools rather than employee API functions.\par

\subsubsection{Query Logging and Evaluation}

\textbf{Purpose.} Record query metrics and measure the quality and safety of answering.\par

\textbf{Operating context.} Query completion writes metrics under the configured privacy setting. An operator runs evaluation under recorded roles and reports outcomes against the pinned dataset and configuration.\par

\Needspace{7\baselineskip}

\textbf{FR-08 Query metrics with controlled content logging}\par

The system shall record query role, retrieved chunk identifiers, relevance scores, latency, and token count; it shall store question and answer text only when content logging is enabled.\par

\textbf{Acceptance.} AT-08: with content logging disabled, required metric fields are present and question/answer text is absent. With it enabled, text may be recorded. Logged role values are bounded and sanitised. The client-deployment default is content logging off.\par

\Needspace{7\baselineskip}

\textbf{FR-09 Repeatable evaluation report}\par

The system shall provide an evaluation script that runs the defined test-question set and reports the resulting scores.\par

\textbf{Acceptance.} AT-09: provide a valid evaluation file, run it as each recorded role, and obtain traceable per-question and aggregate results. Record the dataset revision, configuration, scoring method, and unresolved or failed questions.\par

Evaluation uses bounded concurrency, backoff for provider rate limits, and resumable cached results as specified in R4. Cache reuse must preserve role, source-revision, and configuration boundaries. Numeric settings remain open.\par

\subsubsection{Business Rules}

\begingroup\small

\begin{longtable}{|P{\dimexpr0.0956\linewidth-0.5735\tabcolsep-0.3824\arrayrulewidth\relax}|P{\dimexpr0.7059\linewidth-4.2353\tabcolsep-2.8235\arrayrulewidth\relax}|P{\dimexpr0.1985\linewidth-1.1912\tabcolsep-0.7941\arrayrulewidth\relax}|}

\hline\rowcolor{black!8}

\textbf{ID} & \textbf{Rule} & \textbf{Source} \\ \hline

\endfirsthead

\hline\rowcolor{black!8}

\textbf{ID} & \textbf{Rule} & \textbf{Source} \\ \hline

\endhead

\hline\endfoot

\hline\endlastfoot

BR-01 & The only employee roles are new\_hire, hr, finance, and engineer. Allowed departments are exactly those in Section 2.3. Unknown roles fail closed. & R3 Section 3; R6 config.py \\ \hline

BR-02 & Department assignment follows the longest matching source-path prefix. An explicit file mapping overrides its parent folder. An unmapped source is refused. & R5; R6 departments.py \\ \hline

BR-03 & Paths are relative and use forward slashes. Segments allow ASCII letters, digits, period, underscore, and hyphen. Traversal, empty segments, trailing dots or spaces, and excluded prefixes are refused. & R5 path rules \\ \hline

BR-04 & Access restriction applies to vector and keyword results before final candidate selection and before any passage reaches the model. & R4 Sections 4.1, 4.7 \\ \hline

BR-05 & The not-found threshold uses reranker relevance when enabled, otherwise best vector similarity. A fused RRF score is not the sufficiency measure. Each configuration has a calibrated threshold. & R4 Section 4.9 \\ \hline

BR-06 & Citations identify only chunks supplied for the current generation. Invalid identifiers are removed; no valid citation means not found. & R3 FR-13; R4 Section 4.8 \\ \hline

BR-07 & A changed or deleted document loses its obsolete chunks during offline ingestion. Unchanged documents need not be embedded again. & R4 Section 3.1 \\ \hline

BR-08 & The demo selection contains 190 pages: engineering 60, finance 35, people 50, onboarding 5, and general 40. Each selected page has at least 150 real words by the documented selection method. & R5 selection rules \\ \hline

\end{longtable}

\endgroup

finance/readmes and finance/jihu-support are excluded case-insensitively. Approved mapping matches are case-sensitive; the complete path map remains in R5.\par

\subsection{Performance Requirements}

The following targets are provisional. The measured environment, workload, scoring method, and cost basis must be recorded before a pass is declared.\par

\begingroup\small

\begin{longtable}{|P{\dimexpr0.1544\linewidth-0.9265\tabcolsep-0.6176\arrayrulewidth\relax}|P{\dimexpr0.3971\linewidth-2.3824\tabcolsep-1.5882\arrayrulewidth\relax}|P{\dimexpr0.4485\linewidth-2.6912\tabcolsep-1.7941\arrayrulewidth\relax}|}

\hline\rowcolor{black!8}

\textbf{ID and priority} & \textbf{Requirement and target} & \textbf{Verification and qualification} \\ \hline

\endfirsthead

\hline\rowcolor{black!8}

\textbf{ID and priority} & \textbf{Requirement and target} & \textbf{Verification and qualification} \\ \hline

\endhead

\hline\endfoot

\hline\endlastfoot

NFR-01\newline Medium & The system shall complete at least 95\% of answer requests within 8 seconds under the agreed acceptance workload. & PT-01: measure elapsed time from API receipt of a valid question to the complete answer response. Record hardware, dataset, configuration, concurrent users, request count, and warm-up policy. Concurrency and the final protocol remain OD-01. Source: R3 NFR-01. \\ \hline

NFR-02\newline Medium & The evaluated system shall achieve faithfulness of at least 0.85 and answer correctness of at least 75\% on the agreed test questions. & QT-01: retain per-question scores and aggregate results. The scoring rubric, scoring model if used, treatment of refusals, and denominator must be approved under OD-05 before a pass can be declared. Source: R3 NFR-02. \\ \hline

NFR-04\newline Medium & The average language-model charge shall not exceed USD 0.02 per evaluated question. & CT-01: divide total billed generation cost by the evaluated question count and retain token usage and the price basis used for that run. State whether evaluation-judge calls are excluded. The accounting boundary remains OD-05. Source: R3 NFR-04. \\ \hline

\end{longtable}

\endgroup

The 190-document demo is a reproducible dataset, not a production-capacity limit. Concurrency, throughput, and memory remain OD-01; input-size limits remain OD-04. One API worker does not imply one user.\par

\subsection{Logical Database Requirements}

Each document belongs to one department and has a checksum for its ingested revision. It produces chunks referenced by citations and query records. Retrieval indexes derive from the authoritative chunk collection.\par

\begingroup\small

\begin{longtable}{|P{\dimexpr0.1544\linewidth-0.9265\tabcolsep-0.6176\arrayrulewidth\relax}|P{\dimexpr0.3971\linewidth-2.3824\tabcolsep-1.5882\arrayrulewidth\relax}|P{\dimexpr0.4485\linewidth-2.6912\tabcolsep-1.7941\arrayrulewidth\relax}|}

\hline\rowcolor{black!8}

\textbf{Entity} & \textbf{Required information} & \textbf{Relationship or constraint} \\ \hline

\endfirsthead

\hline\rowcolor{black!8}

\textbf{Entity} & \textbf{Required information} & \textbf{Relationship or constraint} \\ \hline

\endhead

\hline\endfoot

\hline\endlastfoot

Document & document\_id, source\_path, title, department, checksum. & document\_id identifies a source document; a document has one approved department and zero or more derived chunks. \\ \hline

Chunk & chunk\_id, document\_id, department, text, citation title, section or page locator. & chunk\_id is unique. document\_id links to its source. Department agrees with the document; locator supports citation. \\ \hline

Role policy & Role name and permitted department set. & A single access table governs the four employee roles; chunks do not carry a second conflicting role list. \\ \hline

Manifest record & source\_path, department, sha256. & Source paths are unique and sorted. The checksum is a 64-character lowercase hexadecimal SHA-256 value. \\ \hline

Query record & Role, retrieved chunk IDs, relevance scores, final chunk IDs, latency, token count; optional question and answer text. & Identifiers refer to the evidence used by that query. Content fields depend on the content-logging switch. \\ \hline

Evaluation question & id, role, question, expected\_answer, source\_doc, should\_refuse. & The expected answer and refusal flag are written against the pinned source collection. \\ \hline

Evaluation result & Question reference, observed outcome, configuration and dataset revision, reported scores and measurements. & Results remain traceable to the question set, source revision, and scoring conditions. \\ \hline

\end{longtable}

\endgroup

\begingroup\small

\begin{longtable}{|P{\dimexpr0.1029\linewidth-0.6176\tabcolsep-0.4118\arrayrulewidth\relax}|P{\dimexpr0.7426\linewidth-4.4559\tabcolsep-2.9706\arrayrulewidth\relax}|P{\dimexpr0.1544\linewidth-0.9265\tabcolsep-0.6176\arrayrulewidth\relax}|}

\hline\rowcolor{black!8}

\textbf{ID} & \textbf{Integrity or lifecycle rule} & \textbf{Verification} \\ \hline

\endfirsthead

\hline\rowcolor{black!8}

\textbf{ID} & \textbf{Integrity or lifecycle rule} & \textbf{Verification} \\ \hline

\endhead

\hline\endfoot

\hline\endlastfoot

DR-01 & The system shall retain document identity and source metadata needed to trace every served chunk to its source. & AT-02 \\ \hline

DR-02 & The system shall reject chunks labelled with a department outside the approved department set. & AT-02 \\ \hline

DR-03 & The system shall use a single authoritative chunk collection from which retrieval indexes are built or updated. & DT-01 \\ \hline

DR-04 & The system shall remove obsolete chunks for a changed or deleted document from the serving collection during offline ingestion. & DT-01 \\ \hline

DR-05 & The system shall check chunk counts and checksums across the stored collection and derived indexes at startup and refuse to serve when they disagree. & DT-01 \\ \hline

DR-06 & The demo manifest shall use source checksums calculated after CRLF line endings are normalised to LF. & DT-02 \\ \hline

DR-07 & The system shall omit question and answer text from query records when content logging is disabled. & AT-08 \\ \hline

\end{longtable}

\endgroup

DR-01 and DR-02 derive from FR-02 and R4 Section 5; DR-03 to DR-05 from R4 Section 3.1; DR-06 from R5; DR-07 from FR-08. The current ID convention uses the source path for document\_id and a suffix such as \#000 for chunk\_id.\par

Offline ingestion removes obsolete documents and chunks. Log retention, record deletion, and backup policy remain OD-02; credential storage remains OD-03.\par

\subsection{Design Constraints}

R4 selects Python, FastAPI, Streamlit, SQLite, ChromaDB, BM25, local embeddings, a reranker, and Groq. These are design choices, subject to review if the SRS behaviour is preserved. The following operating constraints already affect the baseline.\par

\begingroup\small

\begin{longtable}{|P{\dimexpr0.1029\linewidth-0.6176\tabcolsep-0.4118\arrayrulewidth\relax}|P{\dimexpr0.6176\linewidth-3.7059\tabcolsep-2.4706\arrayrulewidth\relax}|P{\dimexpr0.2794\linewidth-1.6765\tabcolsep-1.1176\arrayrulewidth\relax}|}

\hline\rowcolor{black!8}

\textbf{ID} & \textbf{Constraint} & \textbf{Basis and verification} \\ \hline

\endfirsthead

\hline\rowcolor{black!8}

\textbf{ID} & \textbf{Constraint} & \textbf{Basis and verification} \\ \hline

\endhead

\hline\endfoot

\hline\endlastfoot

DC-01 & Version 1 shall run with a single API worker and perform ingestion offline while serving is stopped. & R4 Sections 3.1 and 8; inspect deployment settings and demonstrate the update procedure. \\ \hline

DC-02 & The backend shall enforce access before evidence reaches generation; client-supplied role choices shall not grant authority. & R4 Sections 4.1 and 4.2; AT-05 and AT-11. \\ \hline

DC-03 & Ingestion and evaluation shall remain operator-run command-line activities without remote document-modification endpoints. & R4 Section 6; inspect routes and operator commands. \\ \hline

DC-04 & The demo shall identify the pinned GitLab Handbook source revision and retain its attribution record. & R5 and R6 README; DT-02 and documentation inspection. \\ \hline

\end{longtable}

\endgroup

\subsubsection{Standards Compliance}

The document follows the supplied course guidance and template. The project records the handbook under CC BY-SA 4.0 and retains its attribution. No regulatory, accessibility, or security certification is claimed.\par

Any domain-specific compliance obligations for a real deployment must be identified before that deployment is approved.\par

\subsection{Software System Attributes}

\subsubsection{Reliability}

\Needspace{7\baselineskip}

\textbf{FR-06 refuses unsupported questions, FR-13 suppresses answers without valid citations, and DR-05 blocks inconsistent serving data. QT-01 evaluates the support for generated claims. No MTBF target has been agreed.}\par

Provider and infrastructure failures are service errors, not evidence that an answer is absent. Recovery and timeout behaviour remain OD-04.\par

\subsubsection{Availability}

The demo depends on its host and model provider. Updates stop and restart the service. No uptime percentage, recovery-time objective, or recovery-point objective is approved; OD-02 and OD-06 cover these decisions.\par

\subsubsection{Security and Privacy}

\begingroup\small

\begin{longtable}{|P{\dimexpr0.1544\linewidth-0.9265\tabcolsep-0.6176\arrayrulewidth\relax}|P{\dimexpr0.4632\linewidth-2.7794\tabcolsep-1.8529\arrayrulewidth\relax}|P{\dimexpr0.3824\linewidth-2.2941\tabcolsep-1.5294\arrayrulewidth\relax}|}

\hline\rowcolor{black!8}

\textbf{ID and priority} & \textbf{Requirement} & \textbf{Verification} \\ \hline

\endfirsthead

\hline\rowcolor{black!8}

\textbf{ID and priority} & \textbf{Requirement} & \textbf{Verification} \\ \hline

\endhead

\hline\endfoot

\hline\endlastfoot

NFR-03\newline High & The system shall disclose no restricted document content across the defined access-control acceptance questions. [R3 NFR-03] & ST-01: inspect answers, citations, retrieved evidence, and model payloads for all role-denial cases. Any leak fails acceptance for the test set. \\ \hline

NFR-07\newline High & No API key, signing secret, password, or other deployment secret shall be committed to the project repository. [R3 NFR-07] & ST-02: inspect tracked files and configuration examples for secrets; confirm sensitive local artifacts are excluded. Record the inspection scope. \\ \hline

\end{longtable}

\endgroup

Access security is specified by FR-05 and FR-10 to FR-13. A public not-found result must not expose a restricted source. Client deployments disable content logging by default; automatic personal-data masking is deferred.\par

The design expires tokens after 30 minutes; immediate revocation is not guaranteed. Credential protection, provisioning, signing, browser token handling, and provider data handling remain OD-03 and OD-04.\par

\subsubsection{Maintainability}

The design uses replaceable model and vector-store interfaces and one role-policy table. A component change must preserve access, refusal, and citation behaviour and pass the affected acceptance checks.\par

No numerical complexity, coverage, or repair-time target is agreed. Maintainability review checks interface separation and continued use of the acceptance procedures.\par

\subsubsection{Portability and Reproducibility}

\begingroup\small

\begin{longtable}{|P{\dimexpr0.1544\linewidth-0.9265\tabcolsep-0.6176\arrayrulewidth\relax}|P{\dimexpr0.4632\linewidth-2.7794\tabcolsep-1.8529\arrayrulewidth\relax}|P{\dimexpr0.3824\linewidth-2.2941\tabcolsep-1.5294\arrayrulewidth\relax}|}

\hline\rowcolor{black!8}

\textbf{ID and priority} & \textbf{Requirement} & \textbf{Verification} \\ \hline

\endfirsthead

\hline\rowcolor{black!8}

\textbf{ID and priority} & \textbf{Requirement} & \textbf{Verification} \\ \hline

\endhead

\hline\endfoot

\hline\endlastfoot

NFR-06\newline Medium & The system shall be set up from its README using Docker in less than 30 minutes under the agreed setup conditions. [R3 NFR-06] & RT-01: use a clean supported host and follow the README. Record host, network, downloads, start/end times, and deviations. Conditions remain OD-06. \\ \hline

NFR-05\newline High & The system shall support Markdown and text-based PDF source documents. Scanned PDFs are not required. [R3 NFR-05] & AT-01: verify both required formats. Record parser limitations and reject or report unsupported content according to the approved ingestion contract. \\ \hline

\end{longtable}

\endgroup

Development uses Windows and PowerShell; tests are intended to remain compatible with Linux. The supported deployment and browser matrix remains OD-06 and OD-07.\par

\subsubsection{Usability Accessibility and Scalability}

IR-01 defines the basic employee interactions. A measured usability target and accessibility conformance level are not approved; OD-07 records the required decisions.\par

Capacity is measured under OD-01. Multi-worker serving, finer permissions, and multiple organisations belong to later releases.\par

\subsection{Organisation of Specific Requirements}

Section 3.2 follows the template's feature organisation. Its alternative outlines by mode, user class, object, stimulus, response, functional hierarchy, and use case are not additional required functions and are not duplicated.\par

\subsection{Additional Comments}

Not found is distinct from invalid authentication, invalid input, and service failure. Citation membership confirms which evidence was supplied; faithfulness evaluation checks whether that evidence supports the answer.\par

\clearpage

\section{Change Management Process}

\subsection{Submitting a Change}

A change request is recorded in a repository issue or equivalent written record. It identifies affected IDs, revised wording, the reason, and expected scope impact. Informal discussion alone does not alter the baseline.\par

\subsection{Impact Assessment}

The owner reviews the effect on permissions, data, interfaces, implementation, verification, cost, and schedule. Access changes require revised denial tests. Performance changes need an agreed protocol or measurements.\par

\subsection{Decision and Approval}

Abdul Samad records acceptance, rejection, or deferral of each request. Changes to an academically approved baseline are submitted to the course reviewer when review is required. This procedure takes effect on document approval.\par

\subsection{Updating the Baseline}

An accepted change updates the requirement, verification, traceability, relevant design, and revision history. Unchanged obligations retain their IDs; withdrawn IDs are not reused. Approved document versions are retained with their implementation and test references.\par

\begingroup\small

\begin{longtable}{|P{\dimexpr0.2500\linewidth-1.0000\tabcolsep-0.7500\arrayrulewidth\relax}|P{\dimexpr0.7500\linewidth-3.0000\tabcolsep-2.2500\arrayrulewidth\relax}|}

\hline\rowcolor{black!8}

\textbf{Change record field} & \textbf{Required entry} \\ \hline

\endfirsthead

\hline\rowcolor{black!8}

\textbf{Change record field} & \textbf{Required entry} \\ \hline

\endhead

\hline\endfoot

\hline\endlastfoot

Change identifier & Unique reference for the request. \\ \hline

Affected requirements & FR, NFR, supporting rule, interface, and acceptance IDs. \\ \hline

Requested change and reason & Current wording, proposed wording, and rationale. \\ \hline

Impact & Scope, data, security, design, tests, cost, and schedule implications. \\ \hline

Decision & Accepted, rejected, or deferred, with the decision maker and date. \\ \hline

Baseline reference & Document version and implementation or test references after approval. \\ \hline

\end{longtable}

\endgroup

Proposed rules and provisional targets are not approval records.\par

\clearpage

\section{Document Approvals}

Approval records agreement on the identified document version. Signatures and dates remain blank until review. Academic approval does not represent enterprise-customer approval.\par

\begingroup\small

\begin{longtable}{|P{\dimexpr0.1912\linewidth-1.9118\tabcolsep-1.1471\arrayrulewidth\relax}|P{\dimexpr0.2500\linewidth-2.5000\tabcolsep-1.5000\arrayrulewidth\relax}|P{\dimexpr0.1250\linewidth-1.2500\tabcolsep-0.7500\arrayrulewidth\relax}|P{\dimexpr0.2500\linewidth-2.5000\tabcolsep-1.5000\arrayrulewidth\relax}|P{\dimexpr0.1838\linewidth-1.8382\tabcolsep-1.1029\arrayrulewidth\relax}|}

\hline\rowcolor{black!8}

\textbf{Role} & \textbf{Name} & \textbf{Decision} & \textbf{Signature} & \textbf{Date} \\ \hline

\endfirsthead

\hline\rowcolor{black!8}

\textbf{Role} & \textbf{Name} & \textbf{Decision} & \textbf{Signature} & \textbf{Date} \\ \hline

\endhead

\hline\endfoot

\hline\endlastfoot

Author and project owner & Abdul Samad & Pending & \_\_\_\_\_\_\_\_\_\_\_\_\_\_\_\_ & \_\_\_\_\_\_\_\_\_\_\_\_ \\ \hline

Course reviewer & Sir Fahim Muhammad Khan & Pending & \_\_\_\_\_\_\_\_\_\_\_\_\_\_\_\_ & \_\_\_\_\_\_\_\_\_\_\_\_ \\ \hline

\end{longtable}

\endgroup

\subsection{Approval Conditions}

Review covers scope, clarity, permissions, sources, and verification. A draft may retain open decisions, but operational acceptance must identify their resolution or explicit deferral.\par

\subsection{Release Acceptance Record}

Product acceptance is recorded separately after testing. The record identifies the product revision, dataset, configuration, environment, procedures, results, unresolved defects, and final decision.\par

\clearpage

\section{Supporting Information}

Sections 6.1 and 6.2 define verification and traceability. Section 6.3 controls open decisions. Sections 6.4 to 6.6 provide provenance, implementation context, and a subject index.\par

\subsection{Verification and Acceptance Plan}

Verification combines tests, inspection, analysis, and demonstration. The core evaluation file, eval/questions.csv, contains the planned 30 questions, with expected answers derived from the pinned source pages.\par

\begingroup\small

\begin{longtable}{|P{\dimexpr0.2574\linewidth-1.5441\tabcolsep-1.0294\arrayrulewidth\relax}|P{\dimexpr0.0809\linewidth-0.4853\tabcolsep-0.3235\arrayrulewidth\relax}|P{\dimexpr0.6618\linewidth-3.9706\tabcolsep-2.6471\arrayrulewidth\relax}|}

\hline\rowcolor{black!8}

\textbf{Question category} & \textbf{Count} & \textbf{Purpose} \\ \hline

\endfirsthead

\hline\rowcolor{black!8}

\textbf{Question category} & \textbf{Count} & \textbf{Purpose} \\ \hline

\endhead

\hline\endfoot

\hline\endlastfoot

Simple factual & 12 & Check direct retrieval, correctness, and citation support. \\ \hline

Multi-section or table & 6 & Check answers requiring evidence from several passages. \\ \hline

Restricted access & 6 & Check refusal and absence of restricted evidence. \\ \hline

Outside the documents & 6 & Check not-found behaviour when no answer exists. \\ \hline

\end{longtable}

\endgroup

Authentication, injection, citation, integrity, and interface fixtures supplement the 30 questions. Permission tests cover every role and allowed/denied department pair, not only the six core refusal questions.\par

\begingroup\small

\begin{longtable}{|P{\dimexpr0.1176\linewidth-0.7059\tabcolsep-0.4706\arrayrulewidth\relax}|P{\dimexpr0.6029\linewidth-3.6176\tabcolsep-2.4118\arrayrulewidth\relax}|P{\dimexpr0.2794\linewidth-1.6765\tabcolsep-1.1176\arrayrulewidth\relax}|}

\hline\rowcolor{black!8}

\textbf{Procedure} & \textbf{Required evidence} & \textbf{Coverage} \\ \hline

\endfirsthead

\hline\rowcolor{black!8}

\textbf{Procedure} & \textbf{Required evidence} & \textbf{Coverage} \\ \hline

\endhead

\hline\endfoot

\hline\endlastfoot

PT-01 & Complete latency samples and percentile calculation under the approved environment and workload. & NFR-01 \\ \hline

QT-01 & Faithfulness and correctness rubric, per-question scores, aggregate calculations, and refusal treatment. & NFR-02 \\ \hline

ST-01 & Zero restricted-content disclosures across the recorded access-control fixtures. & NFR-03 \\ \hline

CT-01 & Tokens, run cost, pricing basis, question count, and mean cost calculation. & NFR-04 \\ \hline

RT-01 & Clean-host setup transcript, conditions, elapsed time, and README deviations. & NFR-06 \\ \hline

ST-02 & Repository-secret inspection and excluded-artifact review with a stated inspection scope. & NFR-07 \\ \hline

DT-01 & Initial load, unchanged update, changed document, deletion, and deliberately inconsistent-index cases. & BR-07; DR-03 to DR-05; DC-01 \\ \hline

DT-02 & Manifest uniqueness, order, quotas, mappings, checksum format, local-file agreement, and attribution record. & BR-08; DR-06; DC-04 \\ \hline

DI-01 & API route inventory and operator-tool demonstration showing no remote ingestion or evaluation route. & DC-03 \\ \hline

\end{longtable}

\endgroup

\textbf{Acceptance rule.} Required procedures must produce their expected evidence. Any restricted-content leak fails safety acceptance. Provisional targets require approved measurement conditions; failures and unexecuted checks remain in the report.\par

\textbf{Comparison control.} Vector-only, hybrid, and hybrid-plus-reranker runs use the same dataset and question set, with their configuration and threshold recorded. Reports include false refusals and false answers.\par

\subsection{Requirements Traceability Matrix}

The matrix links each baseline requirement to its source, SRS section, and verification. Supporting obligations are grouped below it.\par

\begingroup\small

\begin{longtable}{|P{\dimexpr0.1471\linewidth-1.1765\tabcolsep-0.7353\arrayrulewidth\relax}|P{\dimexpr0.4265\linewidth-3.4118\tabcolsep-2.1324\arrayrulewidth\relax}|P{\dimexpr0.1618\linewidth-1.2941\tabcolsep-0.8088\arrayrulewidth\relax}|P{\dimexpr0.2647\linewidth-2.1176\tabcolsep-1.3235\arrayrulewidth\relax}|}

\hline\rowcolor{black!8}

\textbf{ID} & \textbf{Source} & \textbf{SRS section} & \textbf{Verification} \\ \hline

\endfirsthead

\hline\rowcolor{black!8}

\textbf{ID} & \textbf{Source} & \textbf{SRS section} & \textbf{Verification} \\ \hline

\endhead

\hline\endfoot

\hline\endlastfoot

FR-01 & R3 FR-01; R4 3.1 & 3.2.1 & AT-01 \\ \hline

FR-02 & R3 FR-02; R4 5; R5 & 3.2.1 & AT-02 \\ \hline

FR-03 & R3 FR-03; R4 3.2 & 3.2.3 & AT-03, QT-01 \\ \hline

FR-04 & R3 FR-04; R4 5 & 3.2.4 & AT-04 \\ \hline

FR-05 & R3 FR-05; R4 4.1, 4.7 & 3.2.2 & AT-05, ST-01 \\ \hline

FR-06 & R3 FR-06; R4 4.3 & 3.2.3 & AT-06 \\ \hline

FR-07 & R3 FR-07; R4 6 & 3.2.4 & AT-07 \\ \hline

FR-08 & R3 FR-08; R4 4.10, 7 & 3.2.5 & AT-08 \\ \hline

FR-09 & R3 FR-09; R4 7 & 3.2.5 & AT-09 \\ \hline

FR-10 & R3 FR-10; R4 4.6 & 3.2.3 & AT-10 \\ \hline

FR-11 & R3 FR-11; R4 4.2 & 3.2.2 & AT-11 \\ \hline

FR-12 & R3 FR-12; R4 4.4 & 3.2.2 & AT-12 \\ \hline

FR-13 & R3 FR-13; R4 4.8 & 3.2.4 & AT-13 \\ \hline

NFR-01 & R3 NFR-01 & 3.3 & PT-01 \\ \hline

NFR-02 & R3 NFR-02 & 3.3 & QT-01 \\ \hline

NFR-03 & R3 NFR-03 & 3.6.3 & ST-01 \\ \hline

NFR-04 & R3 NFR-04 & 3.3 & CT-01 \\ \hline

NFR-05 & R3 NFR-05 & 3.6.5 & AT-01 \\ \hline

NFR-06 & R3 NFR-06 & 3.6.5 & RT-01 \\ \hline

NFR-07 & R3 NFR-07 & 3.6.3 & ST-02 \\ \hline

\end{longtable}

\endgroup

\begingroup\small

\begin{longtable}{|P{\dimexpr0.1765\linewidth-1.4118\tabcolsep-0.8824\arrayrulewidth\relax}|P{\dimexpr0.4632\linewidth-3.7059\tabcolsep-2.3162\arrayrulewidth\relax}|P{\dimexpr0.1544\linewidth-1.2353\tabcolsep-0.7721\arrayrulewidth\relax}|P{\dimexpr0.2059\linewidth-1.6471\tabcolsep-1.0294\arrayrulewidth\relax}|}

\hline\rowcolor{black!8}

\textbf{Supporting IDs} & \textbf{Source and purpose} & \textbf{SRS section} & \textbf{Verification} \\ \hline

\endfirsthead

\hline\rowcolor{black!8}

\textbf{Supporting IDs} & \textbf{Source and purpose} & \textbf{SRS section} & \textbf{Verification} \\ \hline

\endhead

\hline\endfoot

\hline\endlastfoot

IR-01 to IR-06 & R3 FR-05, FR-07, FR-11, FR-12; R4 Sections 2, 4 and 6 & 3.1 & AT-05, AT-07, AT-11, AT-12 \\ \hline

BR-01 to BR-04 & R3 role policy; R4 access filtering; R5 path rules & 3.2.6 & AT-02, AT-05, AT-11 \\ \hline

BR-05 to BR-08 & R4 thresholds, citations and updates; R5 selection & 3.2.6 & AT-06, AT-13, DT-01, DT-02 \\ \hline

DR-01 to DR-07 & R3 FR-02, FR-08; R4 data integrity; R5 checksums & 3.4 & AT-02, AT-08, DT-01, DT-02 \\ \hline

DC-01 to DC-04 & R4 operating and API constraints; R5 provenance & 3.5 & AT-05, AT-11, DT-01, DT-02, DI-01 \\ \hline

\end{longtable}

\endgroup

\subsection{Open Decisions}

The owner records the decisions below before approving the affected acceptance results. Undecided values do not mean unlimited capacity or indefinite retention.\par

\begingroup\small

\begin{longtable}{|P{\dimexpr0.0956\linewidth-0.7647\tabcolsep-0.4779\arrayrulewidth\relax}|P{\dimexpr0.4706\linewidth-3.7647\tabcolsep-2.3529\arrayrulewidth\relax}|P{\dimexpr0.1838\linewidth-1.4706\tabcolsep-0.9191\arrayrulewidth\relax}|P{\dimexpr0.2500\linewidth-2.0000\tabcolsep-1.2500\arrayrulewidth\relax}|}

\hline\rowcolor{black!8}

\textbf{ID} & \textbf{Decision required} & \textbf{Affected items} & \textbf{Closure evidence} \\ \hline

\endfirsthead

\hline\rowcolor{black!8}

\textbf{ID} & \textbf{Decision required} & \textbf{Affected items} & \textbf{Closure evidence} \\ \hline

\endhead

\hline\endfoot

\hline\endlastfoot

OD-01 & Acceptance hardware; RAM and storage needs; simultaneous users; throughput; latency request count, warm-up, and run conditions. & NFR-01; Sections 2.1.6 and 3.3 & Approved benchmark protocol and first baseline results. \\ \hline

OD-02 & Log and evaluation-record retention; deletion process; backup frequency and location; recovery objectives and procedure. & FR-08; Sections 2.1.7, 3.4, 3.6.2 & Documented operating and retention policy. \\ \hline

OD-03 & Credential schema and storage; account and role provisioning; signing method; provider/model version; private-data handling approval. & FR-11; Sections 3.1.3 and 3.6.3 & Reviewed authentication and deployment configuration specification. \\ \hline

OD-04 & Question and file-size limits; invalid input; corrupt/unsupported files; exact JSON and error schemas; generic reply wording; timeouts; retry bounds; log sanitisation limits. & FR-01, FR-06 to FR-08, FR-12; IR-01 to IR-06 & Approved interface and error contract, with boundary-case tests. \\ \hline

OD-05 & Quality rubric; scoring tool/version; correctness denominator; refusal treatment; calibrated thresholds; cost-accounting scope and price basis. & NFR-02, NFR-04; FR-06, FR-09 & Frozen evaluation protocol and scored baseline report. \\ \hline

OD-06 & Demo hosting target; supported operating systems; Docker prerequisites; whether setup timing includes downloads and credentials. & FR-07; NFR-06 & Deployment record and clean-host setup protocol. \\ \hline

OD-07 & Supported browsers, screen assumptions, keyboard use, accessibility scope, and any usability task threshold. & IR-01; Section 3.6.6 & Agreed interface acceptance checklist. \\ \hline

\end{longtable}

\endgroup

Enterprise stakeholder validation and the signatures in Section 5 remain outstanding.\par

\subsection{Data Provenance and Attribution}

The dataset uses GitLab Handbook commit \nolinkurl{2c77a1f5b8c8a80cb7b5151ff11cfad84f98bfbb}, downloaded on 11 October 2026 according to R5. Selected files are copied unchanged to the local department folders and excluded from Git.\par

\begingroup\small

\begin{longtable}{|P{\dimexpr0.7059\linewidth-2.8235\tabcolsep-2.1176\arrayrulewidth\relax}|P{\dimexpr0.2941\linewidth-1.1765\tabcolsep-0.8824\arrayrulewidth\relax}|}

\hline\rowcolor{black!8}

\textbf{Department} & \textbf{Selected pages} \\ \hline

\endfirsthead

\hline\rowcolor{black!8}

\textbf{Department} & \textbf{Selected pages} \\ \hline

\endhead

\hline\endfoot

\hline\endlastfoot

engineering & 60 \\ \hline

finance & 35 \\ \hline

people & 50 \\ \hline

onboarding & 5 \\ \hline

general & 40 \\ \hline

Total & 190 \\ \hline

\end{longtable}

\endgroup

Each selected page has at least 150 real words after the documented exclusions. Deterministic round-robin selection spreads topics; the manifest is sorted by source path. SHA-256 checksums normalise CRLF to LF.\par

Dataset attribution: GitLab Handbook, GitLab, \url{https://handbook.gitlab.com}, source repository \url{https://gitlab.com/gitlab-com/content-sites/handbook}, licensed as recorded by the project under Creative Commons Attribution-ShareAlike 4.0 International, \url{https://creativecommons.org/licenses/by-sa/4.0/}. The SRS describes the application and does not reproduce handbook policy text.\par

\subsection{Implementation Context}

The reviewed source implements role permissions, department-path mapping, deterministic page selection, checksums, and Markdown chunking with front-matter removal. Title extraction and the complete citation metadata are still in progress.\par

The API, chat, storage, retrieval, generation, logging, evaluation, Docker setup, and live demo remain planned work. Their inclusion in the SRS states the intended release, not completed delivery.\par

\subsection{Subject Index}

\begingroup\small

\begin{longtable}{|P{\dimexpr0.4265\linewidth-1.7059\tabcolsep-1.2794\arrayrulewidth\relax}|P{\dimexpr0.5735\linewidth-2.2941\tabcolsep-1.7206\arrayrulewidth\relax}|}

\hline\rowcolor{black!8}

\textbf{Subject} & \textbf{Location} \\ \hline

\endfirsthead

\hline\rowcolor{black!8}

\textbf{Subject} & \textbf{Location} \\ \hline

\endhead

\hline\endfoot

\hline\endlastfoot

Access permissions and roles & 2.3; 3.2.2; BR-01 and BR-04 \\ \hline

Authentication and tokens & 3.1.3; 3.2.2; 3.6.3 \\ \hline

Citations and source support & 3.2.3 to 3.2.4; AT-04 and AT-13 \\ \hline

Data ingestion and classification & 3.2.1; 3.2.6; 3.4 \\ \hline

Dataset provenance & 1.4; 6.4 \\ \hline

Document approval and revisions & Document Control; Sections 4 and 5 \\ \hline

Evaluation and acceptance & 3.2.5; 6.1 to 6.2 \\ \hline

Logging and privacy & FR-08; DR-07; 3.6.3 \\ \hline

Not-found behaviour & FR-06; FR-12; BR-05; 3.8 \\ \hline

Open decisions & 6.3 \\ \hline

Performance and cost & 3.3 \\ \hline

Portability and setup & 3.6.5 \\ \hline

Release boundary and future scope & 1.2; 2.6 \\ \hline

\end{longtable}

\endgroup
\end{document}
